\documentclass[11pt,a4paper]{article}

% ---------------------------------------------------------------------------
% Classificacao de topicos - case VDE
% Fonte oficial do relatorio. Compila no Overleaf com pdfLaTeX.
% Um unico arquivo: as figuras sao inline (TikZ/pgfplots), sem assets externos.
% ---------------------------------------------------------------------------

\usepackage[T1]{fontenc}
\usepackage[utf8]{inputenc}
\usepackage[brazilian]{babel}
\usepackage[margin=1.9cm]{geometry}
\usepackage{booktabs}
\usepackage{xcolor}
\usepackage{amsmath}
\usepackage{microtype}
\usepackage{enumitem}
\usepackage{titlesec}
\usepackage{tcolorbox}
\usepackage{tikz}
\usetikzlibrary{arrows.meta,positioning,shapes.geometric}
\usepackage{pgfplots}
\pgfplotsset{compat=1.17}
\usepackage[hidelinks]{hyperref}

\definecolor{azul}{HTML}{1F4E79}
\definecolor{cinza}{HTML}{F2F4F7}
\definecolor{verde}{HTML}{2E7D32}
\definecolor{vermelho}{HTML}{B3261E}
\definecolor{ambar}{HTML}{B26A00}

\newcommand{\bd}[1]{\textbf{#1}}
\setlist[itemize]{leftmargin=1.2em,itemsep=1pt,topsep=1pt}
\setlength{\parskip}{0pt}
\titlespacing*{\section}{0pt}{5pt}{2pt}

\tikzset{
  box/.style={draw=azul,rounded corners,fill=cinza,align=center,inner sep=3pt,
              minimum height=8mm,text width=19mm,font=\small},
  arr/.style={-{Latex[length=2mm]},thick,draw=azul}
}

\title{\vspace{-1.4cm}\textbf{\color{azul}Classificação de tópicos: case VDE}\\[3pt]
\large Do raciocínio ao desenho de produção}
\author{Leo Vitor}
\date{}

\begin{document}
\maketitle
\vspace{-1.3cm}

% ============================ SUMARIO EXECUTIVO ============================
\begin{tcolorbox}[colback=cinza,colframe=azul,
                  title=\textbf{Sumário executivo},boxrule=0.6pt,arc=2pt]
\begin{itemize}
\item Operação: 222 publicadas (61,8\% das 359), 136 em revisão, 1 fora do escopo.
\item Qualidade: \bd{0,90 a 0,92} contra dois rotuladores de IA, que entre si concordam em
      0,85. Os rótulos vêm de IA: mede-se concordância, não precisão.
\item Com 60 itens, \bd{não há como separar 90\% de 95\%}: os intervalos se sobrepõem. A meta
      de 95\% permanece \bd{não demonstrada}.
\end{itemize}
\end{tcolorbox}

\noindent\emph{Nota de calibração:} a previsão foi \bd{revisada após a segunda rotulagem}:
precisão $\sim$0,86 (0,74 a 0,94), cobertura $\sim$0,62. O 0,90 central é concordância com um
rotulador de IA, não precisão (o gabarito de especialista não existe aqui).

% ================================= a) COMO PENSEI =========================
\section*{a) Como eu pensei}

O desafio é escolher o tema certo de cada questão entre $\sim$375 temas parecidos, sem mandar
tudo para revisão. O erro nasce em dois pontos, e tratei cada um: (1) o tema certo não está
entre as opções oferecidas ao modelo; (2) o modelo escolhe, mas não sabe quando não tem
certeza. Daí o desenho de duas etapas: uma busca que garante a lista de opções e um juiz com
abstenção, seguido de uma checagem que corrige o segundo ponto.

Uma busca textual propõe candidatos (TF-IDF + pais e filhos). Dois modelos
(\texttt{gemini-3.6-flash} e \texttt{gemini-3.1-flash-lite})
avaliam as opções em quatro votos. Um verificador de 5º voto, com a árvore inteira, entra
\emph{antes} da consolidação e empurra para revisão quando discorda. Só então as regras
consolidam: votos iguais publicam; votos em subtemas do mesmo ramo publicam no nó comum;
ramos diferentes vão para revisão; maioria de \texttt{fora} vira \texttt{fora}.

\begin{figure}[h]
\centering
\begin{tikzpicture}[node distance=6mm]
\node[box] (clean) {Limpeza\\\scriptsize HTML $\to$ texto};
\node[box,right=of clean] (cand) {Candidatos\\\scriptsize TF-IDF + pais/filhos};
\node[box,right=of cand] (juiz) {Juiz\\\scriptsize 2 modelos, 4 votos};
\node[box,below=9mm of juiz] (ver) {Verificador\\\scriptsize 5º voto, árvore inteira};
\node[box,right=of ver] (reg) {Regras\\\scriptsize decide};
\node[box,right=of reg] (out) {publicar / revisar / fora};
\draw[arr] (clean)--(cand);
\draw[arr] (cand)--(juiz);
\draw[arr] (juiz)--(ver);
\draw[arr] (ver)--(reg);
\draw[arr] (reg)--(out);
\end{tikzpicture}
\caption{Fluxo do pipeline entregue. O verificador (5º voto, \emph{antes} das regras) vê a árvore
inteira e pega o erro ``confiante e errado''; não provei que seja a única forma.}
\end{figure}

\bd{Onde eu errei, e o que isso mudou.} Dos 11 erros publicados na base, 6 são de busca (o tema
certo não estava na lista) e 5 de decisão. Todos os 11 foram \textbf{votos unânimes}: quando a
busca traz as opções erradas, os dois modelos concordam no erro, e a concordância sozinha não
sinaliza o problema. Foi por isso que criei o verificador.

\bd{O que testei, e o que descartei.} \emph{Testado e descartado:} política seletiva (95 questões
com $\sim$17 rótulos, na mesma amostra); BM25 e RRF (fusão lexical); embeddings densos
(\texttt{e5-large}: empataram); subir sempre para o pai (+1 caso); SVM linear (0,441). \emph{Descartado por argumento:} cross-encoder, que só reordena a lista e não resolve o
tema que não entra nela.

\bd{O que fiz com o diagnóstico.} Como o erro de busca se concentra em legislação especial,
ataquei esse ponto offline, sem tocar na entrega: um dicionário jurídico (traduz a citação da lei
para o nome que aparece no título) somado à expansão por irmãos. Os resultados estão na Seção~b,
inclusive um teste fim-a-fim que mostrou um limite importante.

% ================================= b) DIAGNOSTICO =========================
\section*{b) Diagnóstico}

\bd{Números (93 questões rotuladas por IA).}

\begin{center}
\begin{tabular}{lr}
\toprule
Métrica & Valor \\
\midrule
Concordância observada (base, 261 publicadas) & 0,838 (57/68) \\
Concordância observada (após verificador, 222) & \bd{0,90 (54/60)} \\
Wilson 95\% da entrega & [0,80; 0,95] \\
Erros publicados na base & 11 (6 de busca, 5 de decisão) \\
\bottomrule
\end{tabular}
\end{center}

\noindent Publicadas ao longo do trabalho, a cada ajuste de regra: 255, 263, 246, 261 e 222.

\bd{Duas ressalvas antes dos números.} (1) Os rótulos vêm de outro LLM (DeepSeek), a partir do
texto da questão e da lista de temas: o 0,90 é concordância entre modelos, não precisão contra
especialistas, que permanece desconhecida.
(2) O recall da lista que o juiz \emph{realmente recebeu} na entrega, já com o fallback, foi de
0,882 (82/93): número inflado (parte das questões recebeu a árvore inteira) e \emph{não}
comparável ao 0,667 da busca pura.

\bd{Como o catálogo atrapalha.} (1) Gêmeos: 8 pares de temas com o mesmo título em disciplinas
diferentes. A régua conta apenas \emph{um} dos dois como certo, e o pipeline não sabe qual.
\emph{O porquê da regra:} publicar no tema errado é o pior erro (o aluno estuda o tópico errado),
enquanto revisar só atrasa; na dúvida entre gêmeos, rebaixar é a aposta segura. Na entrega, a
regra é \emph{publicar o gêmeo quando a disciplina de origem coincide com a raiz do tema}; caso
contrário, rebaixar para revisão. É uma \textbf{premissa adotada} (a organização ainda não
confirmou a convenção): custa \textbf{24 questões rebaixadas} ($\approx$ 6,7 p.p. de cobertura) no
entregue. (2) Lixeiras: temas com muitas questões na origem ($\geq$ 40).
(3) Sobreposição pai/filho, que a estratégia do nó comum explora a favor.

\bd{O que atrapalha nos dados.} (1) A \texttt{disciplina\_origem} coincide com a do tema escolhido
em $\sim$95\% das publicadas rotuladas (3 discordâncias em 60) e serve de desempate; mas 2\% das
questões vêm de fora do recorte (Constitucional, Civil, Tributário, Ambiental) ou sem valor.
(2) O \texttt{assunto\_origem} é enganoso: a Q0026 (Lei Maria da Penha) veio como ``Ação Penal''.
(3) Há questões multitemáticas (Q0016, Q0085), em que o ``tema certo'' é uma escolha, não um fato.

\bd{Bloqueio por conteúdo sensível.} Uma questão (Q0227, cenários de violência sexual em Direito
Penal, com \texttt{assunto\_origem} ``Lesão Corporal'', enganoso) teve a resposta bloqueada pelo
filtro de segurança do modelo (\texttt{PROHIBITED\_CONTENT}). O pipeline trata isso como
abstenção: vai para revisão com o motivo registrado e não repete a chamada. Em produção, um
filtro sem critério jurídico decide o que cai na fila humana, e esse tipo de bloqueio
\emph{pode não ser ajustável} (as proteções contra danos graves costumam não ser configuráveis).
A recomendação é \emph{contar} esses casos e \textbf{mantê-los na revisão humana}, em vez de
afrouxar o filtro. (Um caso em 359 daria $\sim$140 por 50 mil, mas é extrapolação ingênua: o
recorte é só de Penal.)

\bd{Onde o erro se concentra.} Dos 11 erros publicados, \textbf{3 são de legislação especial};
os demais são de Parte Geral, Processo Penal e casos de ambiguidade real. Onde o erro de
\emph{busca} aparece, o juiz escolhe, de forma unânime, o vizinho mais próximo.

\bd{O verificador funciona?} Fixei o critério \emph{antes} de rodar: vetar ao menos 3 erros reais,
com no máximo 20\% de vetos falsos. Ele vetou \bd{5 dos 11 erros} publicados (4 de busca, em que
o tema certo nem fora oferecido) e recusou \bd{3 das 57 certas} (5\%), dentro do critério. Dos 6
erros que escaparam, minha adjudicação (não especialista, em \texttt{docs/adjudicacao\_erros.csv})
diz que 1 favorece o rótulo (Q0120), 1 favorece o modelo (Q0238) e 4 são ambíguos (Q0273, Q0016,
Q0306, Q0085). É leitura de um revisor, não medição.

\bd{Reconciliação do gargalo.} A busca limita a \textbf{precisão das publicadas} (tema certo fora
da lista); agregação, gêmeos e limiar limitam a \textbf{cobertura}. Melhorar só a busca ataca o
primeiro limite, não o segundo (o A/B mostra).

\bd{Ataque ao ponto fraco.} Duas alavancas de busca, testadas offline (93 questões, $k=10$):

\begin{center}
\begin{tabular}{lcc}
\toprule
Mecanismo de busca & \texttt{recall\_hier} & Legislação Especial \\
\midrule
Busca pura (TF-IDF + pais/filhos) & 0,667 & 5/19 \\
+ dicionário (\texttt{top+dic}) & 0,731 & 12/19 \\
+ irmãos, sem dicionário (\texttt{irmao}) & 0,882 & 9/19 \\
\bd{dicionário + irmãos} & \bd{0,957} & \bd{16/19} \\
\bottomrule
\end{tabular}
\end{center}

O dicionário sobe o piso; os irmãos, o topo. No holdout a melhora se mantém (0,613 $\to$ 0,968),
\emph{com a ressalva} de que a última expansão do dicionário foi decidida olhando os erros
remanescentes, um deles no holdout; não é um holdout estritamente cego.

\begin{tcolorbox}[colback=cinza,colframe=ambar,boxrule=0.6pt,arc=2pt,
                  title=\textbf{Teste fim-a-fim (A/B pago, nível da decisão)}]
Levei o dicionário até a decisão na Opus (150 questões, 148 com tema; quase fora da amostra, 41 já
no conjunto de desenho), mesmos modelos e regras, com o verificador:
\begin{center}
\begin{tabular}{lcc}
\toprule
Configuração (limpa) & \texttt{recall\_hier} & Publicação correta (c/ verificador) \\
\midrule
\texttt{filho} (controle) & 0,791 & 0,886 (78/88) \\
\texttt{filho+dic} (dicionário) & 0,845 & 0,888 (79/89) \\
\bottomrule
\end{tabular}
\end{center}
O dicionário sobe o recall (+5,4 p.p.) mas muda a decisão em +0,2 p.p. (1 questão). A maior
parte dos acertos extras cai em \texttt{revisar}: o juiz acha o tema, mas a agregação não
publica, e a métrica conta só as publicadas. No pareado: +2 conserta, $-1$ quebra, 8 neutras.
Não foi adotado na entrega.
\end{tcolorbox}

\bd{Duas rotulagens, lado a lado.} A concordância do pipeline é medida contra dois rotuladores
de IA; nenhum é o gabarito de especialistas, e os dois não são independentes entre si (41 questões
em comum):

\begin{center}
\begin{tabular}{lccc}
\toprule
Rotulagem & Itens & Concordância & Wilson 95\% \\
\midrule
DeepSeek (conjunto de desenho das regras) & 60 & 0,90 & 0,80 a 0,95 \\
Opus (checagem quase fora da amostra) & 88 & 0,92 & 0,84 a 0,96 \\
Entre os dois rotuladores (régua pai/filho) & 41 & 0,85 & 0,72 a 0,93 \\
\bottomrule
\end{tabular}
\end{center}

O resultado central (0,90) vem do conjunto que orientou o desenho das regras; a checagem com a
Opus (0,92) repete 41 dessas questões, então é \emph{quase} fora da amostra. O achado central: as
duas rotulagens concordam no ramo, não no tema exato (0,54 no exato e 0,85 na régua pai/filho nas
41 comuns). Cuidado com os dois ``88'': o \bd{81/88} é o pipeline entregue na Opus; o
\textbf{78/88} é o braço \texttt{filho} do A/B acima, uma configuração diferente, não o mesmo
conjunto.

\bd{Quanto custa precisar mais.} Como 95\% é a pergunta central, medi o trade-off entre
cobertura e precisão (varredura de limiar no conjunto de validação). Os intervalos de Wilson
abaixo mostram por que \emph{não} dá para cravar 95\%:

\begin{center}
\begin{tabular}{lccc}
\toprule
Cobertura & Precisão observada & $n$ & Wilson 95\% \\
\midrule
0,73 & 0,84 & 32 & [0,68; 0,93] \\
\bd{0,62} & 0,90 (entrega, c/ verificador) & 60 & [0,80; 0,95] \\
0,59 & 0,86 & 28 & [0,69; 0,94] \\
0,43 & 0,93 & 15 & [0,70; 0,99] \\
0,26 & 1,00 & 6 & [0,61; 1,00] \\
\bottomrule
\end{tabular}
\end{center}

\begin{figure}[h]
\centering
\begin{tikzpicture}
\begin{axis}[
  width=0.74\textwidth, height=2.7cm,
  xlabel={Cobertura (fração publicada)}, ylabel={Precisão observada},
  xmin=0.2, xmax=0.8, ymin=0.55, ymax=1.05,
  grid=both, minor tick num=1,
  legend style={font=\footnotesize, at={(0.02,0.02)}, anchor=south west, fill=white},
  tick label style={font=\footnotesize}, label style={font=\small}
]
\addplot+[mark=*,thick,color=azul,
  error bars/.cd, y dir=both, y explicit]
coordinates {
 (0.724,0.844) += (0,0.086) -= (0,0.164)
 (0.708,0.844) += (0,0.087) -= (0,0.166)
 (0.677,0.871) += (0,0.077) -= (0,0.151)
 (0.646,0.867) += (0,0.080) -= (0,0.157)
 (0.593,0.857) += (0,0.083) -= (0,0.167)
 (0.471,0.857) += (0,0.100) -= (0,0.200)
 (0.432,0.933) += (0,0.057) -= (0,0.233)
 (0.265,1.000) += (0,0.000) -= (0,0.390)
};
\addlegendentry{base (varredura de limiar)}
\addplot[only marks,mark=star,mark size=3.2pt,color=verde] coordinates {(0.618,0.900)};
\addlegendentry{entrega (c/ verificador)}
\addplot[dashed,thick,color=vermelho] coordinates {(0.2,0.95)(0.8,0.95)};
\addlegendentry{meta 95\%}
\end{axis}
\end{tikzpicture}
\caption{Cobertura $\times$ precisão, com intervalos de Wilson. Os intervalos se sobrepõem: com
esta amostra não dá para separar 90\% de 95\%. A cobertura alta só aparece com $n \leq 15$ e
intervalo largo.}
\end{figure}

% ================================= c) PRODUCAO ============================
\section*{c) Produção (50 mil questões por semana)}

Esta seção é proposta; não foi executada nesta entrega (orçamento limitado de API e prazo curto).

\bd{Encaixe no fluxo.} (1) as questões novas entram pelo importador; (2) a busca propõe
candidatos; (3) um classificador leve resolve a maioria dos casos com probabilidade calibrada;
(4) o LLM julga o que sobra e um verificador checa os ambíguos; (5) o que ainda diverge vai para
a fila de revisão humana; (6) os reportes de alunos e a revisão alimentam a amostra auditada que
mede a qualidade semana a semana.

\bd{Arquitetura, do barato para o caro.} (1) Recuperação ampliada: TF-IDF + dicionário jurídico
e irmãos (grátis); embeddings jurídicos densos só com GPU. (2) Classificador treinado: \emph{só
depois de acumular rótulos dos professores} (hoje não existe; o SVM deu 0,441). (3) LLM com
verificação seletiva e revisão humana nas divergências.

\bd{Custo por etapa.} Medido onde há registro (fallback e A/B, em \texttt{uso.json} e
\texttt{uso\_etapas.jsonl}); no juiz base, estimado; no verificador, lido no console
(246 chamadas; os tokens não foram gravados em arquivo):

\begin{center}
\begin{tabular}{lcc}
\toprule
Etapa (359 questões) & Tokens de entrada & Custo (US\$, preços reais) \\
\midrule
Juiz (4 votos) & $\sim$2,1 M & $\sim$1,04 \\
Fallback (28 questões $\times$ 4) & $\sim$0,9 M (medido) & $\sim$0,47 \\
Verificador (246 chamadas) & $\sim$2,14 M (lido no console) & $\sim$0,56 \\
\midrule
\bd{Total} & \bd{$\sim$5,2 M} & \bd{$\sim$2,1} \\
\bottomrule
\end{tabular}
\end{center}

Preços: \texttt{gemini-3.6-flash} (0,75 / 3,75 por milhão) e \texttt{gemini-3.1-flash-lite}
(0,25 / 1,50). \bd{$\sim$14,4 M tokens de entrada por 1.000 questões} (não os 33,8 M que supunham
a árvore inteira em toda questão). O verificador tem \emph{mais} tokens por chamada, mas custa
\emph{menos} por questão que os quatro votos do juiz, porque usa só o modelo barato. É um
\bd{piso}, só tokens de entrada (sem raciocínio nem saída). Custo por questão \bd{entre
$\sim$US\$ 0,006} (piso) \bd{e $\sim$R\$ 0,09} (teto observado: R\$ 32,98 das rodadas
$\div$ 359, com os experimentos); para 50 mil por semana, de $\sim$US\$ 290 a $\sim$R\$ 4,6 mil.

\bd{Cache de contexto.} O bloco repetido (instruções + árvore, $\sim$9 mil tokens) pode ser
enviado uma vez e reaproveitado. O plano em tokens é grátis (\texttt{-{}-dry-run}: $\sim$72\% de
economia na entrada); o teste real é \emph{pequeno}, a \textbf{saída não foi guardada} e os
créditos de API acabaram, então não o reexecutei. Reproduzir com
\texttt{python -m src.contexto\_cache -{}-criar} (\texttt{docs/cache\_contexto.md}); o desconto real
em R\$ só se confirma no contrato.

\bd{Limitações honestas.} (1) Juiz e verificador são da mesma família (Gemini): erram juntos, e a
Q0120 passou pelos dois. (2) O \texttt{THRESH} e as lixeiras quase não filtram (confianças perto
de 0,95; nas lixeiras o erro foi \emph{menor}, 0,12 contra 0,21). (3) As regras (verificador,
gêmeos) foram desenhadas na amostra de 60 em que são medidas. (4) A verificação encarece com a
árvore inteira por chamada.

\bd{Medição contínua.} Três métricas semanais, além do tamanho da amostra: (i) precisão na
amostra auditada; (ii) taxa de correção da fila (quantos revisados estavam errados); (iii)
reportes de alunos por tema, como sinal de erro que a auditoria não pegou. Para separar 90\% de
95\% são necessários \bd{$\sim$434 itens por grupo} (poder 80\%); em torno de 0,90, $\sim$139 para
$\pm$5 p.p. A amostra auditada precisa ser da ordem de \textbf{centenas}, não de dezenas.

\bd{O que eu faria com mais tempo.} (1) Um verificador de \emph{outra família} (hoje juiz e
verificador são ambos Gemini e podem errar juntos); (2) rótulos de especialistas nos 20 a 30
desacordos; (3) irmãos só nos casos de baixo risco; (4) o cache em produção; (5) confirmar a
convenção dos gêmeos. \textbf{Com o professor no circuito:} (6) ligações curadas lei$\to$tema e
entre gêmeos, validadas por ele; (7) decisão por custo, com probabilidade calibrada; (8)
aprendizado ativo: a revisão dele gera rótulos e treina o classificador.

\bd{Meta de 95\%.} Sem gabarito humano não há garantia formal: não afirmamos que os 95\% foram
atingidos.

% ================================= d) MEMO ================================
\clearpage
\section*{(d) Memo para o professor}
{\small\itshape Este memo é para você, especialista em Direito, não em IA. Uma página, sem
jargão.}

\vspace{4pt}
\bd{O que eu construí}
\begin{itemize}
\item Um sistema que decide, para cada uma das 359 questões, entre publicar, revisar ou fora.
\item Resultado: 222 publicadas (61,8\%), 136 em revisão e 1 fora do escopo.
\item 0,90 a 0,92 de concordância contra duas avaliações feitas por IA (que entre si concordam em
      0,85), mas nenhuma é o gabarito humano: a meta de 95\% não está confirmada.
\end{itemize}

\bd{O que encontrei}
\begin{itemize}
\item Dos 11 erros publicados, só 3 são de legislação especial; alguns vêm de questões com
      assertivas I--IV, em que o enunciado dá pouca pista e o tema certo nem chega a ser oferecido
      ao modelo. É um problema de busca, não de conhecimento jurídico.
\item O erro é confiante: nos casos errados, os dois modelos concordaram por unanimidade. Por
      isso adicionei a verificação, um 5º julgamento que vê todos os temas e manda para
      revisão quando discorda.
\item O catálogo tem armadilhas: temas com o mesmo nome em disciplinas diferentes e temas que se
      sobrepõem (pai e filho dizendo o mesmo). \emph{Não observei} erro por escolha de gêmeo,
      mas são poucos casos e a convenção ainda não foi confirmada.
\end{itemize}

\bd{O que já melhorei (offline, sem tocar na entrega)}
\begin{itemize}
\item Um dicionário jurídico somado à busca por temas irmãos elevou o achado do tema (ou de um
      relacionado) de 0,667 para 0,957; em legislação especial, de 5/19 para 16/19. A parte
      local é grátis.
\item Testei isso até o fim: o ganho de busca não se converteu em publicação (+0,2 ponto
      percentual, 1 questão). Por isso não mudei a entrega.
\end{itemize}

\bd{O que recomendo mudar nos temas}
\begin{itemize}
\item Dar nomes que distingam os gêmeos (ex.: citar a disciplina no título). São 8 pares, com
      ids: 4684/6716, 5063/5466, 4688/6717, 4702/6723, 4704/6724, 4703/6734, 4705/6725,
      4708/6728.
\item Fundir ou hierarquizar melhor subtemas que repetem o conteúdo do pai.
\item Marcar temas ``guarda-chuva'' (onde aceitar o nível geral é seguro) e os que exigem o
      nível específico.
\end{itemize}

\bd{O que preciso de você}
\begin{itemize}
\item Uma amostra rotulada por você, começando por \textbf{20 a 30 casos de desacordo} entre as
      duas avaliações de IA, para medir a precisão de verdade. Hoje a referência é assistida por
      IA e só mede concordância.
\item Confirmação da convenção dos temas gêmeos no gabarito.
\item Um parecer sobre 6 casos em que o modelo e o rótulo de IA discordam (primeiro o que o
      modelo publicou, depois o que o rótulo sugere):
      \begin{itemize}
      \item Q0120 (provas; resposta à acusação): Teoria Geral das Provas $\times$ Defesa Preliminar.
      \item Q0238 (reforço da fiança, CPP): Liberdade Provisória $\times$ Prisão em Flagrante.
      \item Q0273 (crime punido na modalidade culposa): Crime Doloso e Culposo $\times$ Condescendência Criminosa.
      \item Q0016 (jurisprudência sobre prisão): Teoria Geral das Prisões $\times$ Prisão Preventiva.
      \item Q0306 (associação criminosa): Crimes Contra à Paz Pública $\times$ Organizações Criminosas.
      \item Q0085 (concussão, agente funcionário público): Concurso de Pessoas $\times$ Concussão e Excesso de Exação.
      \end{itemize}
\end{itemize}

\end{document}
